% concatenated from Templated_Order_of_H.tex
%%%%%%%%%%%%%%%%%%%%%%%%%%%%%%%%%%%%%%%%%%%%%%%%%%%%%%%
% % Lines starting with % are comments, which are ignored.
% % This is a handy way of indicating the date and version of
% % your document, to wit:
% %
% % LaTeX sample file
% % Modified March, 2002
% %
%%%%%%%%%%%%%%%%%%%%%%%%%%%%%%%%%%%%%%%%%%%%%%%%%%%%%%%
% % Title and author(s)
%%%%%%%%%%%%%%%%%%%%%%%%%%%%%%%%%%%%%%%%%%%%%%%%%%%%%%%
\title{The Order of the Non-universal Tree of a Hilbert Space with Respect to the Haar Basis}
\author{Sam Whitmire
	 \thanks{Indiana University Bloomington Department of Mathematics, Bloomington, IN 47405}
                  }
%%%%%%%%%%%%%%%%%%%%%%%%%%%%%%%%%%%%%%%%%%%%%%%%%%%%%%%
\documentclass{article}
%%%%%%%%%%%%%%%%%%%%%%%%%%%%%%%%%%%%%%%%%%%%%%%%%%%%%%%
% %
% % The next command allows your in import encapsulated
% % postscript files, .epsf or .eps files, which
% % contain vector graphic image data.
% %
%%%%%%%%%%%%%%%%%%%%%%%%%%%%%%%%%%%%%%%%%%%%%%%%%%%%%%%
\usepackage{graphicx}
\usepackage{amssymb}
\usepackage{amsmath}
\usepackage{mathtools}
\usepackage{amsthm}
\usepackage{geometry}
\usepackage{biblatex}
\addbibresource{dst.bib}
%%%%%%%%%%%%%%%%%%%%%%%%%%%%%%%%%%%%%%%%%%%%%%%%%%%%%%%
% % We use newtheorem to define theorem-like structures
% %
% % Here are some common ones. . .
%%%%%%%%%%%%%%%%%%%%%%%%%%%%%%%%%%%%%%%%%%%%%%%%%%%%%%%
\newtheorem{theorem}{Theorem}
\newtheorem{lemma}[theorem]{Lemma}
\newtheorem{proposition}[theorem]{Proposition}
\newtheorem{scolium}[theorem]{Scolium}   %% And a not so common one.
\newtheorem{definition}[theorem]{Definition}
\newtheorem{corollary}[theorem]{Corollary}
\newenvironment{AMS}{}{}
\newenvironment{keywords}{}{}
\newcommand{\im}{\text{im}}
\newcommand{\spn}{\text{span}}
\newcommand{\N}{\mathbb N}
\newcommand{\R}{\mathbb R}
\newcommand{\C}{\mathbb C}
\newcommand{\floor}[1]{\lfloor #1 \rfloor}
\DeclarePairedDelimiter{\ceil}{\lceil}{\rceil}
%%%%%%%%%%%%%%%%%%%%%%%%%%%%%%%%%%%%%%%%%%%%%%%%%%%%%%%
% %   The first thanks indicates your affiliation
% %
% %  Just the name here.
% %
% % Your mailing address goes at the end.
% %
% % \thanks is also how you indicate grant support
% %
%%%%%%%%%%%%%%%%%%%%%%%%%%%%%%%%%%%%%%%%%%%%%%%%%%%%%%%


\begin{document}
\newpage
\maketitle
%%%%%%%%%%%%%%%%%%%%%%%%%%%
% abstract, keywords and Subject classification are optional.
%%%%%%%%%%%%%%%%%%%%%%%%%%%
\begin{abstract}
In \cite{Bo}, Bossard introduced the notion of the non-universal tree $T_{NU}(X)$ associated with each separable Banach space $X$ which does not contain an isomorphic copy of $C(2^\N)$. Together with the order operation defined on well-founded trees, we obtain a method of classifying the complexity of separable Banach spaces by the degree of isomorphism of finite-dimensional subspaces of $C(2^\N)$. Despite further refinement of this concept in later years, we are unaware of any specific instances of direct computations of the order of the non-universal tree for a concrete space and basis. We show that if $H$ is any separable Hilbert space, then $o(T_{NU}(H)) = \omega + 1$ when taken with respect to the Haar basis for $C(2^\N)$, demonstrating that the class of Hilbert spaces is the least complex class with respect to this measurement when considering this basis.
\end{abstract}

% Most people don't use these, so they are "commented out"
% by starting the lines with a "%"
%\begin{keywords}
%   \LaTeX, typesetting
%\end{keywords}

%\begin{AMS}
%   50C60, 18C25
%\end{AMS}

%%%%%%%%%%%%%%%%%%%%%%
% % Here is the start of the Text
%%%%%%%%%%%%%%%%%%%%%%
\section{Introduction}
The conceptualization of the construction of the NU tree for a space $X$ can be traced back to Bourgain in \cite{Bu}. In \cite{Bo}, Bossard formalized the development and demonstrated that $X\mapsto o(T_{NC}(X))$ is a $\boldsymbol{\Pi}_1^1$-rank on $NU$, the collection of all separable Banach spaces which do not contain an isomorphic copy of $C(2^\N)$. At each step of the construction, one asks for an isomorphic equivalence between the spaces spanned by finite sequences in $X$ and finite lists of elements of a Schauder basis of $C(2^\N)$; hence this order operation measures the degree to which $X$ contains isomorphic finite-dimensional subspaces of $C(2^\N)$ despite not containing an isomorphic copy of $C(2^\N)$ itself.

Much interest has been expressed regarding this mapping. For example, in \cite{AD}, the question of dependence of $o(T_{NU}(X))$ on the choice of basis used in the construction was tackled, and they were able to verify that there is a $\boldsymbol\Pi_1^1$ rank on NU which uniformly bounds $o(T_{NU}(X,(e_n)))$ for each normalized Schauder basis $(e_n)$ of $C(2^\N)$. Even so, to our knowledge, there are no concrete examples of actual computations of $o(T_{NU}(X,(e_n)))$ for any space $X$ and any basis $(e_n)$ of $C(2^\N)$. This paper aims to fill that gap with the following result:

\begin{theorem}\label{MR}
Let $H$ be a separable Hilbert space and $(f_{t_n})_{n=0}^\infty$ be the Haar basis for $C(2^\N)$. Then $\\o(T_{NC}(X,(f_{t_n}))) = \omega + 1$.
\end{theorem}

We also prove that $o(T_{NC}(X))\geq\omega + 1$ for any infinite-dimensional $X$, so this result indicates that Hilbert spaces lie at the absolute minimum of the complexity spectrum of Banach spaces which do not contain $C(2^\N)$ and for which comparison is made with the Haar basis. In certain cases, we are even able to prove direct estimates on the lengths of the vectors contained in the tree, and we derive these in the pentultimate section of the paper.

\section{Preliminaries}
In this section, we present a rapid review of the fundamental concepts required to discuss the construction of the tree in Theorem \ref{MR}. For a more detailed exposition, we invite the reader to consult \cite{K}. 

To construct the desired objects, we utilize a series of ideas united from functional analysis and descriptive set theory. By $2^\N$ we mean the set of all binary sequences under the metric 
\begin{equation*}
d(\alpha,\beta) = \begin{cases}
\frac{1}{n+1} & n \text{ is the least integer for which $\alpha(n)\neq\beta(n)$}\\
0 & \alpha = \beta
\end{cases}
\end{equation*}
It is well-known that $2^\N$ is a compact zero-dimensional Polish space under this metric. By $C(2^\N)$ we mean the set of all real-valued continuous functions on $2^\N$ equipped with the usual supremum norm
\begin{equation*}
||f||_\infty = \sup_{\alpha\in 2^\N} \{|f(\alpha)|\}
\end{equation*}
which makes $C(2^\N)$ into a Banach space. This space has a natural Schauder basis called the Haar basis which will stand at the center of our result. Let $2^{<\N}$ be the set of all finite binary sequences and let $\mathcal O\subseteq 2^{<\N}$ consist of the empty string together with all binary strings which end in 0. For $s,t\in \mathcal O$, define $s\cap t$ to be the longest string which is an initial segment of both $s$ and $t$. Also define $s \prec t$ if and only if $|s| < |t|$ or [$|s| = |t|$ and at the first index $i$ for which $s(i)\neq t(i)$, we have $s(i) = 0$ and $t(i) = 1$].

With a bit of effort, we have the following proposition (pp. 109-110 in \cite{D}):

\begin{proposition}
There is a unique bijection $h : \mathcal O\rightarrow\N$ which satisfies the following properties:
\begin{enumerate}
\item If $|s| < |t|$, then $h(s) < h(t)$.
\item If $|s|=|t|$ and $s \prec t$, then $h(s) < h(t)$.
\end{enumerate}
\end{proposition}

By setting $t_n = h^{-1}(n)$, we obtain an enumeration of the strings in $\mathcal O$. Here is a sample listing of the first few values of $h$:
\begin{align*}
&h(\varnothing) = 0 \\
& h((0)) = 1 \\
&h((00)) = 2 \qquad  h((10)) = 3 \\ 
&h((000)) = 4 \qquad h((010))=5 \qquad h((100)) = 6 \qquad h((110))=7 \qquad \\
&\ldots
\end{align*}
Now let $V_n = \{\alpha\in 2^\N\mid \alpha|_{|t_n|} = t_n\}$. This is a clopen subset of $2^\N$, and so its characteristic function $\chi_{V_n} : 2^\N\rightarrow\R$ is continuous. We set $f_{t_n} = \chi_{V_n}$. These are linearly independent, and by the Stone-Weierstrass theorem, $\overline{\spn\{f_{t_0}, f_{t_1},\ldots\}} = C(2^\N)$. Hence $(f_{t_n})$ forms a normalized Schauder basis of $C(2^\N)$. Moreover, this basis can be shown to be monotone (Claim 6.25 in \cite{D}).

Let $A$ be a countable set. A tree over $A$ is a set $T\subseteq 2^{<A}$ which is downards closed, i.e. if $t\in T$ and $s$ is a substring of $t$, then $s\in T$ as well. $T$ is said to be ill-founded if there is a sequence $\alpha = (a_n)\in A^\N$ with $\alpha|_k\in T$ for every $k\in\N$, and is said to be well-founded otherwise. Let $Tr_A = \{T\subseteq A^{\N}\mid T\text{ is a tree}\}$ and $WF_A = \{T\in Tr_A\mid T\text{ is well-founded}\}$. For a well-founded tree $T$, we can define the iterated derivatives $\{T^{(\xi)}\mid \xi<\omega_1\}$  of $T$ by the rules
\begin{align*}
&T^{(0)} = T\\
&T^{(\xi+1)} = \{t\in T\mid (\exists s\in T)[s\sqsupset t]\}\\
&T^{(\lambda)} = \bigcap_{\xi < \lambda} T^{(\xi)} \text{ when $\lambda$ is a limit ordinal}
\end{align*}
We think of the successor derivative operation as trimming the furthest leaves off of each branch of our tree, while the limit derivative operation consolidates all of these trimmings together. When $T$ is well-founded, there must exist $\xi < \omega_1$ with $T^{(\xi)} = \varnothing$; the smallest such $\xi$ is called the order of the tree and we write $o(T) = \xi$. There is also a phrasing of this operation in terms of well-founded relations, from which it follows that the order of every well-founded tree is a successor ordinal (cf. pp. 10-11 in \cite{K}).

For a Polish space $X$ and $W\subseteq X$, $W$ is called analytic (written $W\in\boldsymbol\Sigma_1^1(X)$) if $W$ is the image of $\N^\N$ under a continuous function. $W$ is called coanalytic (written $W\in\boldsymbol\Pi_1^1(X)$) if $X\setminus W$ is analytic. A function $\varphi:C\rightarrow\omega_1$ on $C\in\boldsymbol\Pi_1^1(X)$ is called a $\boldsymbol\Pi_1^1$-rank on $X$ if there are relations $\leq_\Sigma\in \boldsymbol\Sigma_1^1(X\times X), \leq_\Pi\in\boldsymbol\Pi_1^1(X\times X)$ such that for each $y\in B$,
\begin{equation*}
\{x\in B\mid \varphi(x)\leq\varphi(y)\} = \{x\in B\mid x\leq_\Sigma y\} = \{x\in B\mid x\leq_\Pi y\}
\end{equation*}

The following result is standard and can be found on any text on descriptive set theory, e.g. Theorem 27.1 in \cite{K}.

\begin{proposition}
$WF_A\in\boldsymbol\Pi_1^1(Tr_A)$, and for $T\in WF_A$, $T\mapsto o(T)$ is a $\boldsymbol\Pi_1^1$ rank on $WF_A$.
\end{proposition}

Now we discuss the construction of the trees. Let $X$ be a separable Banach space.  Fix $C\in[1,\infty)$ and a separable Banach space $Y$. Given $m\in\N$, two sequences $(x_n)_{n=0}^{m-1}\subseteq X, (y_n)_{n=0}^{m-1}\subseteq Y$ are said to be $C$-equivalent if for each sequence $(a_n)_{n=0}^{m-1}\subseteq\R$, we have
\begin{equation*}
\frac{1}{C}\bigg|\bigg|\sum_{n=0}^{m-1} a_n x_n\bigg|\bigg|\leq \bigg|\bigg|\sum_{n=0}^{m-1} a_n y_n\bigg|\bigg|\leq C\bigg|\bigg|\sum_{n=0}^{m-1} a_n x_n\bigg|\bigg|
\end{equation*}
Throughout, we will assume that the space $Y$ of interest for the above definition is $C(2^\N)$, and we will also assume that the sequence $(y_n)$ is an initial segment of the Haar basis $(f_{t_n})_{n\in\N}$ for $C(2^\N)$ defined above. Define the tree $\textbf{T}_{NC}(X,C(2^\N),(f_{t_n}),C)$ on $X$ by the rule $(x_n)_{n=0}^{m-1}\in \textbf{T}_{NC}(X,C(2^\N),C)$ if and only if $(x_n)_{n=0}^{m-1}$ is $C$-equivalent to $(f_{t_n})_{n=0}^{m-1}$. The following fact due to Bourgain is straightforward to verify (see, for instance, \cite{D}, Lemma 2.16):

\begin{theorem}\label{WFR}
$X$ does not contain a copy of $C(2^\N)$ if and only if for every $C\in [1,\infty)$, $\\ \textbf{T}_{NC}(X,C(2^\N),(f_{t_n}),C)$ is well-founded.
\end{theorem}

The theory of well-founded trees is substantially more developed for $\N$ than it is for Banach spaces, so we seek to convert this into a tree on $\N$ rather than $X$. Moreover, to apply this theory, we need this transformation to occur in a Borel manner. To do so, we utilize the universal space for separable Banach spaces. It is well known that if $Z$ is any separable Banach space, then $Z$ is isometrically isomorphic to a closed subspace of $C(2^\N)$. When endowed with the Effros Borel structure inherited from the set of all closed \textit{subsets} of $C(2^\N)$, the collection of all closed \textit{subspaces} of $C(2^\N)$ becomes a standard Borel space which contains an isometric copy of every separable Banach space. We denote this standard Borel space by SB. (This construction lies at the heart of a great number of advancements in functional analysis; more information about the basic facts can be found in Chapters 1 and 2 of \cite{D}.) By the Kuratowski-Ryll-Nardzewski selection theorem (see e.g. \cite{K}, Theorem 12.13), there are Borel functions $\{d_n : \text{SB}\setminus\{\varnothing\}\rightarrow C(2^\N)\}$ such that for every $X\in \text{SB}\setminus\{\varnothing\}$, $\{d_n(X)\}_{n\in\N}$ is dense in $X$. Using these, we define a tree $T_{NU}(X,C(2^\N),(f_{t_n}),k)$ ($k\in\N$ with $k\geq 1$) on $\N$ by the rule $(\ell_0,\ldots,\ell_{m-1})\in T_{NU}(X,C(2^\N),(f_{t_n}),k)$ if and only if $(d_{\ell_n}(X))_{n=0}^{m-1}$ is $k$-equivalent to $(f_{t_n})_{n=0}^{m-1}$. It is easily verifiable that an analogue to Theorem \ref{WFR} holds for this new countable collection of trees. Finally, to condense them into one tree, we define the tree $T_{NU}(X) = T_{NU}(X,C(2^\N), (f_{t_n}))$ on $\N$ by the rule $(k,\ell_0,\ldots, \ell_{m-1})\in T_{NU}(X)$ if and only if $(\ell_0,\ldots, \ell_{m-1})\in T_{NU}(X,C(2^\N), (f_{t_n}), k)$. Note that vacuously we have $\varnothing\in T$ and $(k)\in T$ for every $k\in\N$ with $k\geq 1$, which ensures that $T$ is genuinely a tree. This is a Borel construction in the sense that the map $X\mapsto T_{NC}(X, C(2^\N), (f_{t_n}))$ is a Borel map (following the argument presented in Lemma 2.4 in \cite{D}), and since $\text{WF}$ is a $\boldsymbol\Pi_1^1$ subset of $\text{Tr}$, we have the following result due to Bossard (recorded in \cite{D} as Theorem 2.17):

\begin{theorem}\label{WFN}
The set $\text{NU} = \{X\in\text{SB}\mid X\text{ does not contain an isomorphic copy of $C(2^\N)$}\}$ is a $\boldsymbol\Pi_1^1$ subset of SB, and the map $X\mapsto o(T_{NU}(X))$ is a $\boldsymbol\Pi_1^1$-rank on NU.
\end{theorem}

\section{Main Result}

To prove Theorem \ref{MR}, our first step is an obvious reduction. The definitions of each of the classes of trees in the construction of $T_{NC}(H)$ all follow rules related to satisfaction of inequalities by arbitrary linear combinations of vectors. As a result, if $X$ and $Y$ are isometrically isomorphic, then $o(T_{NU}(X,(f_{t_n})) = o(T_{NU}(Y,(f_{t_n})))$. This allows us to assume without loss of generality that $H = \ell^2$. 

We start with the lower bound in the theorem. In general, one cannot say much about the order of the non-universal tree of an arbitrary separable Banach space beyond the following simple estimate, but in our case, it turns out to be sufficient.

\begin{theorem}
Let $X\in \text{NU}$ be infinite-dimensional. Then $o(T_{NU}(X))\geq\omega + 1$.
\end{theorem}

\begin{proof}
Since $X$ is infinite-dimensional, $X$ must have an infinite linearly independent subset of vectors, and by density, we can find such a family using the selection functions $\{d_n\}$. Extracting to a subsequence if necessary, let $(d_n(X))_{n=0}^\infty$ be linearly independent. Fix $m$ and consider the finite sequences $(d_n(X))_{n=0}^{m-1}$ and $(f_{t_n})_{n=0}^{m-1}$. Let $T : \overline{\spn\{d_0(X),\ldots, d_{m-1}(X)\}}\rightarrow \overline{\spn\{f_{t_0},\ldots, f_{t_{m-1}}\}}$ be defined by $T(d_n(X)) = f_{t_n}$. This is a linear operator between finite dimensional spaces and hence must be continuous. Also, $T$ is invertible, so $T^{-1}$ is also continuous. Hence there exists a number $k_m$ such that for every $(a_n)_{n=0}^{m-1}\subseteq \C$, we have
\begin{equation*}
\frac{1}{k_m}\bigg|\bigg|\sum_{n=0}^{m-1} a_n d_n(X)\bigg|\bigg|\leq \bigg|\bigg|\sum_{n=0}^{m-1} a_n f_{t_n}\bigg|\bigg|\leq k_m\bigg|\bigg|\sum_{n=0}^{m-1} a_n d_n(X)\bigg|\bigg|
\end{equation*}
which implies that $(k_m, 0, \ldots, m-1)\in T_{NU}(X)$ and $o(T_{NU}(X))\geq m+1$. Since this holds for every $m$, $o(T_{NU}(X))\geq\omega$. Finally, as mentioned earlier, $o(T)$ must be a successor ordinal, and so $o(T_{NU}(X))\geq\omega + 1$.
\end{proof}

To find the corresponding upper bound, we introduce a new piece of terminology.

\begin{definition}
Let $T$ be a tree on a set $A$. An \textbf{infinite hand} in $T$ is a sequence $(s_n)_{n=0}^\infty\subseteq T$ with the following properties:
\begin{enumerate}
\item $(|s_n|)_{n=0}^\infty\subseteq\N$ is strictly monotone increasing.
\item There exists an element $s\in T$ with $s\sqsubset s_n$ for every $n$; the element of largest length which satisfies this property is called the \textbf{palm} of the infinite hand.
\end{enumerate}
\end{definition}

\begin{lemma}
Let $T$ be a well-founded tree on a set $A$ with $o(T)\geq\omega+1$. Then $o(T) = \omega+1$ if and only if the only palm of an infinite hand in $T$ is $\varnothing$.
\end{lemma}

\begin{proof}
If $o(T) = \omega+1$ and $T$ contains an infinite hand with palm $s\neq\varnothing$, then $s\in T^\omega = \bigcap_{n\in\omega} T^n$ so that $o(T) > \omega+1$, which is a contradiction. Conversely, if $o(T) > \omega+1$, then there is an element $t\in T^{\omega+1}$ and hence an element $s\in T^\omega\setminus\{\varnothing\}$; without loss of generality, we may assume $s$ is a string that has maximum length among all strings in $T^\omega$. Fix $n$. Since $s\in T^n$, there is an extension $s_0^n$ of $s$ in $T^{n-1}$; then we can find an extension $s_1^n\in T^{n-2}$ of $s_0^n$; iterating this procedure produces a string $s_n\in T^0 = T$ with $|s_n| = |s| + n$ and $s\sqsubset s_n$. The family $(s_n)_{n=0}^\infty$ is an infinite hand in $T$ with palm $s$.
\end{proof}

Theorem \ref{MR} follows immediately from the following lemma by the definition of $T_{NU}(\ell^2, (f_{t_n}))$ as an amalgamation of the various $\textbf{T}_{NU}(\ell^2, (f_{t_n}), k)$ for $k\geq 1$.

\begin{lemma}\label{ML}
Let $C\geq 1$. There exists a number $N\in\N$ depending on $C$ such that if $t\in\textbf{T}_{NU}(\ell^2, (f_{t_n}),C)$, then $|t|\leq N$. That is, there is no infinite hand in $\textbf{T}_{NU}(\ell^2, (f_{t_n}),C)$.
\end{lemma}

\begin{proof}
For notational convenience, let $\textbf{T}(C) = \textbf{T}_{NU}(\ell^2, (f_{t_n}), C)$. Suppose that $(t_i)_{i=0}^\infty$ is an infinite hand in $\textbf{T}(C)$. Write $t_i = (v_0^i,\ldots, v_{m_i-1}^i)$. Without loss of generality, we may assume $m_i = i+1$ for every $i$ (i.e. $|t_i|=i+1$ for every $i$). We will produce the quantity $N$ by looking at a specific subsequence of the $v_n^i$ vectors.

By assumption, $t_i\in \textbf{T}(C)$ for every $i$. This provides us with the inequalities
\begin{equation}\label{r1}
\frac{1}{C}\bigg|\bigg|\sum_{n=0}^{i} a_n f_{t_n}\bigg|\bigg|\leq \bigg|\bigg|\sum_{n=0}^{i} a_n v_n^i\bigg|\bigg|\leq C\bigg|\bigg|\sum_{n=0}^{i} a_n f_{t_n}\bigg|\bigg|
\end{equation}
for each sequence $(a_n)_{n=0}^i\in \R^{i+1}$. Observe that for $n\in\{2^k-1\mid k\in\N\setminus\{0\}\}$, the characteristic functions $\{f_{t_n}\}$ have disjoint support. Hence
\begin{equation*}
\bigg|\bigg|\sum_{n\in\{2^k-1\}, n\leq i} a_n f_{t_n}\bigg|\bigg|_{\infty} = \sup_n\{|a_n|\} = ||(a_n)||_\infty
\end{equation*}
For each $i$, assume the coefficients $(a_n)_{n=0}^i$ are zero unless $n = 2^k - 1$ for some $k\geq 1$. For simplicity, let us relabel the vectors $\{v_n^i\mid n = 2^k - 1, n\leq i\}$ as $(v_k^i)_{k=1}^{p_i}$, where $p_i = \floor{\log_2(i+1)}$. Also, let $\{u_k^i\}$ be an orthonormal basis for $\spn\{v_k^i\}$. There is a unitary transformation sending $u_k^i$ to $e_k$ in $\ell^2_{p_i}$, and in this reduction, the images of the $v_k^i$ are obtained via a change of basis operator $T$. Combining these together, inequality \ref{r1} becomes
\begin{equation}\label{r2}
\frac{1}{C}||(a_k)_{k=1}^{p_i}||_\infty\leq \bigg|\bigg|T\bigg(\sum_{k=1}^{p_i} a_k e_k\bigg)\bigg|\bigg|\leq C||(a_k)_{k=1}^{p_i}||_\infty
\end{equation}
This illustrates that $T$ can be regarded as an isomorphism in $B(\ell_{p_i}^2, \ell_{p_i}^\infty)$ with $||T||\leq C$ and $||T^{-1}||\leq C$, so that $||T||\cdot||T^{-1}||\leq C^2$.

We now apply a theorem recorded by Tomczak and Jaegermann (Proposition 37.6 in \cite{TJ}) regarding the Banach-Mazur distance of the spaces $\ell^2_m$ and $\ell^\infty_m$, which states that
\begin{equation*}
\inf\{||S||\cdot||S^{-1}||\mid S \in B(\ell^2_m,\ell^\infty_m) \text{ is an isomorphism}\} = \sqrt{m}
\end{equation*}
Hence we have
\begin{equation*}
\sqrt{p_i} = \inf\{||S||\cdot||S^{-1}||\mid S \in B(\ell^2_{p_i},\ell^\infty_{p_i}) \text{ is an isomorphism}\} \leq ||T||\cdot||T^{-1}|| \leq C^2
\end{equation*}
Now $p_i$ is increasing and tends to infinity as $i\rightarrow\infty$, which contradicts this inequality. Therefore inequality \ref{r1} can only be satisfied when $p_i\leq C^4$; taking $N = 2^{\ceil{C^4}-1}$ provides the stated upper bound.
\end{proof}

\section{Improved Estimates}

This section contains some stronger estimates on the lengths of the strings of vectors in $\ell^2$ under certain additional conditions. Our starting point is the map $S_m : \ell^2_m\rightarrow \spn\{f_{t_0},\ldots, f_{t_{m-1}}\}$ ($m\geq 1$) defined by $S(e_n) = f_{t_n}$. We would like to determine the operator norm of $S_m$. Observe that
\begin{equation*}
\bigg|\bigg|\sum_{n=0}^{m-1} a_n f_{t_n}\bigg|\bigg|_\infty = \sup_{\alpha\in 2^\N}\bigg\{\bigg|\sum_{n=0}^{m-1} a_n \chi_{V_n}(\alpha)\bigg|\bigg\} = \sup_{F\in\mathcal F_m}\bigg\{\bigg|\sum_{n\in F} a_n\bigg|\bigg\}
\end{equation*}
for some finite collection $\mathcal F_m\subseteq\mathcal P([0..m-1])$ depending on $m$. Estimating the operator norm of $S_m$ can therefore be done by estimating the operator norm of an operator of the form $T_\mathcal{F} : \ell^p_m\rightarrow\ell^{\infty}_{b_m}$ of the form $T_\mathcal{F}(a_0,\ldots, a_{m-1}) = (\sum_{n\in F_0} a_n, \ldots, \sum_{n\in F_{b_m}} a_n)$, where $b_m = |\mathcal F|$.

\begin{lemma}\label{ptoinf}
Let $\mathcal F\subseteq\mathcal P([0..m-1])$ be nonempty and define $T_\mathcal{F}\in B(\ell_m^p,\ell_{b_m}^\infty)$ by 
\begin{equation*}
T_\mathcal{F}(a_0,\ldots, a_{m-1}) = \bigg(\sum_{n\in F_0} a_n, \ldots, \sum_{n\in F_{b_m}} a_n\bigg)
\end{equation*} 
Let $k = \sup_{F\in\mathcal{F}}\{|F|\}$ and $q$ be the Holder conjugate of $p$. Then $||T_\mathcal{F}|| = k^\frac{1}{q}$.
\end{lemma}

\begin{proof}
By Holder's inequality, we have
\begin{align*}
||T_\mathcal{F}((a_0,\ldots, a_{m-1}))||_\infty = \sup_{F\in\mathcal F}\bigg\{\bigg|\sum_{n\in F} a_n\bigg|\bigg\}\leq &\sup_{F\in\mathcal F}\bigg\{\sum_{n\in F}| a_n|\bigg\}\leq \sup_{F\in\mathcal F}\{|F|^\frac{1}{q}\:||a||_p\}= k^\frac{1}{q}||a||_p
\end{align*}
Sharpness of the estimate comes from choosing $G\in\mathcal F$ with $|G| = \sup_{F\in\mathcal F}\{|F|\}$ and taking $a = \chi_G\in\ell^p_m$.
\end{proof}

\begin{corollary}\label{basicvectorestimate}
Let $S_m$ be as above and $m\geq 1$. Then $||S_m|| = \sqrt{2+\floor{\log_2(m-1)}}$.
\end{corollary}

\begin{proof}
By Lemma \ref{ptoinf}, we can determine the value of $||S_m||$ by finding $\sup_{F\in\mathcal F_m}\{|F|\}$ in terms of $m$ for a certain choice of $\mathcal F_m$. For each $a = (a_0,\ldots, a_{m-1})$, $S_m(a)$ is of the form $\sum_{n=0}^{m-1} a_n\chi_{V_n}$, and two sets $V_{n}$, $V_{n'}$ intersect if and only if one is a subset of the other by definition of the basis of $C(2^\N)$. Hence any $S_m(a)$ must be the sum over a set $F\subseteq[0..m-1]$ for which there is an $\alpha\in 2^\N$ such that $\alpha\in V_n$ for every $n\in F$, or equivalently, $t_{|F|}\sqsubseteq \alpha$ (where $t_n$ is determined by the enumeration function $h$). If we let $\mathcal F_m$ be the collection of all such $F$, then by definition of $h$, we have
\begin{equation*}
\sup_{F\in \mathcal F_m}\{|F|\} = 1+\sup\{k\mid 2^{k-1}\leq m-1\} = 2 + \floor{\log_2(m-1)}
\end{equation*}
which gives us the desired result.
\end{proof}

This result tells us that if $m > 1+2^{C^2-2}$, then $(e_0,\ldots, e_{m-1})\notin \textbf{T}(C)$, which gives a stronger bound than that of Lemma \ref{ML} when we consider the standard basis.  Unsurprisingly, since we can unitarily transform the standard basis of $\ell^2$ into any other orthonormal basis and the defining condition of $\textbf{T}(C)$ is invariant under unitary transformations, we have the following result.

\begin{corollary}\label{reorder}
Let $\textbf{T}(C)$ be defined as above and $(e_n)$ be the standard basis of $\ell^2$. Fix $m\geq 1$ and let $(u_n)_{n=0}^{m-1}$ be an orthonormal set in $\ell^2$. Then $(u_n)_{n=0}^{m-1}\in \textbf{T}(C)$ if and only if $(e_n)_{n=0}^{m-1}\in \textbf{T}(C)$. 

 In particular, every string $s\in \textbf{T}(C)$ consisting entirely of orthonormal vectors satisfies $|s|\leq 2^{C^2-2}$.
\end{corollary}

By loosening our grasp on the uniform bound, we can obtain an analogous improved estimate for strings consisting of pairwise-orthogonal vectors (not necessarily orthonormal).

\begin{lemma}\label{boundedorthogonal}
Suppose that $(u_0,\ldots, u_{m-1})\in \textbf{T}(C)$ is a family of pairwise-orthogonal vectors. Set $A = \inf_{n<m}\{||u_n||\}$, $B = \sup_{n < m} \{||u_n||\}$, and $C' = \max\{C B, \frac{C}{A}\}$. Then $m \leq 2^{(C')^2-2}$. 

In particular, for every sting $s\in \textbf{T}(C)$ consisting of pairwise-orthogonal vectors, $|s|\leq 2^{C^4-2}$.
\end{lemma}

\begin{proof}
Without loss of generality (Lemma \ref{reorder}), we may assume that $u_n = b_n e_n$ for some $b_n\in\mathbb R\setminus\{0\}$ (so that $b_n = ||u_n||$ for each $n$). Then we have the inequalities
\begin{equation*}
A\bigg|\bigg|\sum_{n=0}^{m-1} a_n e_n\bigg|\bigg| \leq \bigg|\bigg|\sum_{n=0}^{m-1} a_n b_n e_n\bigg|\bigg| \leq B\bigg|\bigg|\sum_{n=0}^{m-1} a_n e_n\bigg|\bigg|
\end{equation*}
for every $(a_n)_{n=0}^{m-1}\subseteq \R$. Combining this with the fact that $(u_0,\ldots, u_{m-1})\in \textbf{T}(C)$ by assumption, we obtain further inequalities
\begin{equation*}
\frac{1}{C B} \leq \bigg|\bigg|\sum_{n=0}^{m-1} a_n f_{t_n}\bigg|\bigg| \leq \bigg|\bigg|\sum_{n=0}^{m-1} a_n e_n\bigg|\bigg| \leq \frac{C}{A} \bigg|\bigg| \sum_{n=0}^{m-1} a_n f_{t_n}\bigg|\bigg|
\end{equation*}
which implies that
\begin{equation*}
\frac{1}{C'} \leq \bigg|\bigg|\sum_{n=0}^{m-1} a_n f_{t_n}\bigg|\bigg| \leq \bigg|\bigg|\sum_{n=0}^{m-1} a_n e_n\bigg|\bigg| \leq C' \bigg|\bigg| \sum_{n=0}^{m-1} a_n f_{t_n}\bigg|\bigg|
\end{equation*}
By Lemma \ref{reorder}, $m\leq 2^{(C')^2-2}$, giving us the desired result.

For the last statement, let $s = (u_0,\ldots, u_{m-1})\in\textbf{T}(C)$. By definition, this means that 
\begin{equation*}
\frac{1}{C} \leq \bigg|\bigg|\sum_{n=0}^{m-1} a_n f_{t_n}\bigg|\bigg| \leq \bigg|\bigg|\sum_{n=0}^{m-1} a_n u_n\bigg|\bigg| \leq C \bigg|\bigg| \sum_{n=0}^{m-1} a_n f_{t_n}\bigg|\bigg|
\end{equation*}
for every $a_0,\ldots, a_{m-1}\in \mathbb R$. Given $n$, take $a_\ell = 0$ for $\ell\neq n$ and $a_n = 1$. We obtain
\begin{equation*}
\frac{1}{C} = \frac{1}{C} \bigg|\bigg|\sum_{n=0}^{m-1} a_n f_{t_n}\bigg|\bigg| \leq \bigg|\bigg|\sum_{n=0}^{m-1} a_n u_n\bigg|\bigg| \leq C \bigg|\bigg| \sum_{n=0}^{m-1} a_n f_{t_n}\bigg|\bigg| = C
\end{equation*}
or
\begin{equation*}
\frac{1}{C}\leq ||u_n|| \leq C
\end{equation*}
Hence $\frac{1}{C}\leq A \leq B\leq C$ so that $C^2\geq \max\{C B,\frac{C}{A}\}$, which proves the claim.
\end{proof}

\section{Conclusion}

There is still much to be discovered in this direction. The most obvious natural question is to ask about the order of other spaces with respect to the basis $(f_{t_n})$. Natural choices include $\ell^p$ and $L^p([0,1])$ for $p\in(1,\infty)$. Since the duals of these spaces are separable, none of these spaces can contain $C(2^\N)$, and so their $NU$ trees (with respect to $(f_{t_n})$) must have countable order. The proposition from \cite{TJ} cited in the proof of Lemma \ref{ML} suggests that for the spaces $\ell^p$ ($p\in (1,\infty)$) we should have $o(T_{NU}(\ell^p, (f_{t_n}))) = \omega + 1$ as well, but we do not have a proof of this at the moment. 

When working with these trees, the choice of Schauder basis of $C(2^\N)$ is not typically a life or death decision, but there are some associated subtleties. In \cite{PS}, Pelczy\'{n}ski and Singer illustrated that for every Banach space $X$ with a Schauder basis, there are uncountably many normalized Schauder bases of $X$ which are not equivalent in the sense described in the Preliminaries. Argyros and Dodos proved that there is a $\boldsymbol \Pi_1^1$-rank $\varphi$ on $NU$ which satisfies $\varphi(X)\geq o(T_{NU}(X,(e_n)))$ for every $X\in NU$ and Schauder basis $(e_n)$ of $C(2^\N)$ (cf. Theorem 2.27 in \cite{D}; originally proved in \cite{AD}). This implies that $o(T_{NU}(H, (e_n)))$ is bounded above by some countable ordinal $\xi$ regardless of the choice of the basis $(e_n)$ of $C(2^\N)$. However, our work does not provide an estimate as to what this upper bound should be. In other words, it is possible that for a different choice of Schauder basis $(e_n)$ of $C(2^\N)$ which is not equivalent to $(f_{t_n})$, we may have $o(T_{NU}(H,(e_n))) > \omega + 1$. We are interested to know if this is the case.

\printbibliography

\end{document}
